\section{Ameaças à validade e limitações}
\label{sec:ameacas}

\subsection{Validade de construção}

\begin{itemize}[leftmargin=*]
  \item \textbf{$S(c)$ conta ataques, não a sua gravidade.} Com os pesos usados,
    todos os ataques valem o mesmo; a exclusão por caso de uso compensa em parte,
    porque uma propriedade obrigatória em falta não é trocada por custo.
  \item \textbf{A escuta no \emph{broker} mede um subscritor estranho.} Na matriz
    de ataques, quem escuta é um cliente sem credenciais, que K1 já impede. Um
    \emph{broker} curioso ou comprometido lê o conteúdo em K1, K2 e K9, que não
    cifram perante o \emph{broker}, e $S(c)$ não o mostra. Quem o mostra é o
    filtro por caso de uso (a telemetria sensível exige o modo de cifra) e o
    cenário que compara o que um subscritor lê com TLS e com o envelope.
  \item \textbf{Um ataque travado no \emph{broker} conta como impedido.} As
    configurações com transporte protegido impedem muitos ataques por o atacante
    não chegar ao \emph{broker}; isso deixa de valer se o \emph{broker} for o
    adversário. K10 e os cenários com \emph{broker} ativo medem o que resta sem
    essa barreira.
  \item \textbf{\emph{Broker} ativo só em MQTT sem TLS.} O intermediário não
    termina TLS; um \emph{broker} malicioso real fá-lo-ia. O que o envelope deteta
    não depende do transporte, pelo que o resultado se mantém, mas não foi medido
    com TLS.
\end{itemize}

\subsection{Validade interna}

\begin{itemize}[leftmargin=*]
  \item \textbf{O instrumento pesa na medição.} O tempo de processador do
    \emph{gateway} inclui a contabilidade das medições, e a sua memória inclui as
    amostras de latência da janela e a tabela de mensagens vistas, que cresce nos
    primeiros minutos de cada laboratório. As diferenças entre configurações são
    válidas; os valores absolutos de memória são um majorante.
  \item \textbf{Tempo de ligação.} Varia em degraus de cerca de 100~ms por esperas
    internas da biblioteca cliente; usa-se o tempo de processador por ligação.
  \item \textbf{Relógios.} A latência entre máquinas depende da correção do
    relógio dos sensores, com incerteza igual a metade do melhor tempo de ida e
    volta observado, no máximo 0,21~ms na campanha (0,18~ms nas execuções de
    custo). É da ordem das diferenças de latência mediana entre configurações
    vizinhas (K0 122~$\mu$s, K8 179~$\mu$s), mas a mesma para todas: como a ordem é
    aleatória, desloca as medidas sem favorecer nenhuma configuração. A deriva é
    corrigida (Secção~\ref{sec:metodo}); uma alteração do ritmo do relógio a meio
    de uma execução não o é, e por isso o NTP esteve desligado durante a campanha.
    O desvio entre as duas máquinas cresceu de forma regular até 289~ms (cerca
    de 16~ppm), sempre corrigido.
  \item \textbf{Temperatura e frequência.} Nos três pilotos não houve
    \emph{throttling} (\texttt{throttled} sempre \texttt{0x0}), a frequência
    esteve sempre em 2\,400~MHz e a temperatura entre 44 e 52~°C. O mesmo na
    campanha: \texttt{0x0} e 2\,400~MHz em todas as execuções, temperatura entre
    44,6 e 52,4~°C (o máximo na subida de carga). A ordem aleatória distribui o
    efeito que houver pelas configurações.
  \item \textbf{Sensores simulados.} Os sensores são processos Python num
    computador, não dispositivos restritos: o custo no emissor não foi medido no
    nível~2 (o nível~1 mede envio e receção).
  \item \textbf{Uma só máquina de carga.} Os dois agentes de sensores e o atacante
    correm no mesmo computador. Ele gerou 80\,000~mensagens/s nos pilotos sem
    ficar aquém da taxa pedida, mas o que acontece nele (escalonamento, recolha de
    memória do Python) entra na latência medida.
  \item \textbf{Picos de latência.} Em todas as configurações, incluindo K0, que
    não tem criptografia, surgem nos pilotos picos de 22 a 40~ms no p99,9 e no
    máximo, a cargas intermédias. Por existirem em K0, vêm do ambiente e não das
    camadas medidas; não afetam o p99 abaixo de cerca de 300~mensagens/s por
    sensor. Na campanha, os maiores valores nas execuções de custo foram de
    0,5~ms em K0. A partir de 400~mensagens/s por sensor
    (20\,000~mensagens/s), a latência mediana sobe de cerca de 0,2 para 7--14~ms
    em todas as configurações, sem perda, com o \emph{gateway} entre 40\,\% (K0)
    e 63\,\% (K8) de processador, longe da saturação, e volta a descer quando a
    carga aumenta. O salto acontece à mesma taxa por sensor em todas as
    configurações; foi registado, mas a sua origem não foi determinada. Afeta a
    latência acima de
    $R_{\text{alvo}}$ (200~mensagens/s por sensor) e fica abaixo de $L_{\max}$,
    pelo que não muda nenhum veredicto, mas a latência na subida de carga acima
    desse ponto é um majorante.
  \item \textbf{Serviços alheios.} O Docker esteve ativo no Raspberry~Pi durante o
    teste de fumo e os pilotos, e foi parado para a campanha, tal como o ambiente
    gráfico; os pilotos só serviram para fixar parâmetros, não entram nos
    resultados.
  \item \textbf{O próprio laboratório tinha defeitos.} Três foram encontrados nas
    duas primeiras campanhas completas e corrigidos antes da terceira, a única
    usada (Secção~\ref{sec:desvios}). O mais grave fazia contar como impedidos
    ataques que não chegaram a ser feitos; foi detetado por o ataque de inundação
    aparecer impedido em K7, que não tem limitação. Na campanha usada,
    verificou-se em todas as células dos ataques internos que o atacante se ligou
    como o dispositivo comprometido e enviou mensagens.
\end{itemize}

\subsection{Validade externa}

\begin{itemize}[leftmargin=*]
  \item Um só modelo de \emph{hardware} restrito (Raspberry~Pi~5), um só
    \emph{broker} (Mosquitto) e uma rede local. Um microcontrolador sem sistema
    operativo, ou uma rede com perdas, podem ordenar as configurações de outra
    forma.
  \item \textbf{Sem AES em \emph{hardware}} foi emulado desligando extensões no
    OpenSSL do mesmo processador, o que não reproduz a memória nem a
    \emph{cache} de um processador mais simples.
  \item Carga sintética: no nível~2, mensagens de 16~octetos de conteúdo a taxa
    constante, QoS~0, 50 sensores; o tamanho só foi variado no nível~1. Com
    mensagens maiores, o peso do envelope e a vantagem do AES em \emph{hardware}
    crescem (Secção~\ref{sec:aceleracao}).
\end{itemize}

\subsection{Limitações do protótipo}

As limitações da \PUm{Secção~6.11} mantêm-se, com uma alteração: o reinício do
recetor deixou de reabrir a janela anti-repetição (Secção~\ref{sec:prototipo}), à
custa de exigir uma rotação de chaves depois de cada reinício. Continuam por
resolver a distribuição de chaves fora do laboratório, o armazenamento de chaves
(só em memória), a supressão e o atraso de mensagens por um \emph{broker} ativo, e
a confidencialidade e a prova de origem em simultâneo (cifra seguida de
assinatura).
